\scanpage{98}
où $m$ est le nombre des $f_i$ , et $mp$ est bien indépendant de $n$.

La difficulté qui \deleted{reste} semble rester est de trouver les suites exactes que tu demandes de telle façon qu'elles se recollent, pour $n$ variable. Donc, $u$ étant déjà choisi, de trouver un homomorphisme $A^p\quad A^q$ \deleted{donnant $\unsure{u|n}$} dont l'image soit $\operatorname{Ker}u$ ... Si tu y arrives, on pourrait présenter encore 10.10.5 sous forme de trois conditions équivalentes, mais en demandant dans b) une ``prsentation finie\deleted{''} uniformément en $n$'' . Sinon, il faudrait trouver un contre-exemple à \deleted{l'équiva} l'implication a)\quad c) , car il me semble que la question devrait être tirée au clair. En tout état de cause, il faudra donner,en corollaire,l'équivalence des conditions suivantes, valables sans hypothèses noethériennes:

a) $F$ est localement libre de type fini.

b) $F$ est i\unsure{\deleted{i}s|\deleted{s}i}omorphe à la limite projective d'une suite $(F_n)$ de $\unsure{\text{\uuline{$O$}}|\text{\uline{$O$}}}_{X_n}$-Modules loc libres de type fini qui se recollent.

c) $F$ est isomorphe à un $M$ , avec $M$ un $A$-module projectif de type fini.

Pour la démonstration, on procède comme dans 10.10.5 en ut\uline{i}lisant EGA IV 18.3.2.1 . Ce lemme devrait d'ailleurs venir en corollaire \deleted{\unsure{7.|à|0}} après $\unsure{0|O}_I$ 7.2.9.

Pour les modifications que je préconaisais pour I 10.11 , elles tombent à l'eau s\unsure{\deleted{o}i|\deleted{i}o} on n'arrive pas à arranger 10.10.5 sans hypothèses noethériennes; on peut cependant dire que si $F$ sur $X$ est de présentation finie, alors il est limite projective de faisceaux $F_n$ de présentation finie sur les $X_n$ qui se recollent (mais on n'a pas une réciproque), et que $F$ est localement libre s\unsure{\deleted{s}i|\deleted{i}s}s les $F_n$ le sont. \unsure{\deleted{\unsure{S}}E|\deleted{E}\unsure{S}}t \deleted{10.11.4} les autres énoncés du n\unsure{0|\deleted{0}°|\deleted{°}0|°} 10.11. restent valables\unsure{\deleted{,}|x} sans hypothèse noe\unsure{\deleted{h}t|\deleted{t}h}hérienne, sauf 10.7.11.2 et la partie ``injectivité'' dans 10.11.9 (sauf erreur).
% archive sheet number: 97
